\chapter{Summary and Outlook} 
\label{chap:summary_outlook}

\section{Summary} 
\label{sec:summary_findings} 
\paragraph{Problem.}
This thesis addressed the absence of a systematic approach for verifying the security functions of industrial gateways at the physical IT/OT interface. With Regulation (EU) 2023/1230, cybersecurity became a mandatory part of the conformity assessment of machinery, yet the established functional and safety validation process did not verify security functions in a requirement-driven manner.

\paragraph{Approach.}
The methodology was developed with the Design Science Research Methodology. It derives the required tests along a traceable chain. It maps the obligations of Regulation (EU) 2023/1230 through the draft standard prEN 50742 to the component requirements of IEC 62443-4-2, links each requirement to the affected assets, to a STRIDE threat, and to the security function that treats it, and specifies each test case on two levels, a product-neutral scenario and a concrete execution. It reuses the company test framework through a tool-to-target applicability assessment and classifies each test with the verification taxonomy of IEC 62443-4-1 Practice SVV. The methodology was demonstrated on a real Safe Remote I/O device in an isolated laboratory under the direct-access network position NV1.

\paragraph{Main findings.}
The derivation produced fourteen security requirements from Annex III of Regulation (EU) 2023/1230 and mapped each of them to prEN 50742 and to IEC 62443-4-2. The device was decomposed into nine security-relevant assets. The applicability assessment classified the ten company tests as five applicable, four not applicable, and one blocked: network segmentation, default credentials, TLS configuration, and hardening validation were not applicable, because the embedded device exposes no second network zone, no authentication endpoint, no TLS endpoint, and no operating-system shell, and the known-CVE scan was blocked in the laboratory. Of the three network access variants, only NV1 was realised. Fifteen test activities were at least partially executed under NV1, and their evidence was stored. Two requirements were assessed as FAIL on the basis of an evidenced device deviation: the intervention log carries only a relative system tick and no absolute timestamp, which fails the time-attribution obligation of RQ-003, and its content is cleared by a cold restart, which fails the retention obligation of RQ-012. A structured fuzzing run additionally surfaced a temporary, self-recovering disruption of the HTTP service that requires follow-up by development. The remaining requirements were assessed as PARTIAL, inconclusive, or not assessed, which reflects the bounded scope of the practical phase.

\paragraph{Research questions.}
RQ1 is answered by the traceable derivation chain and the two-level catalogue, which turn regulatory obligations into concrete, requirement-linked tests along the chain from requirement to asset, threat, security function, and test case. RQ2 is answered by the applicability assessment, which identified the non-applicable and blocked tests before execution and thereby prevented four tools from being run against interfaces that the device does not expose. RQ3 is answered in part. NV1 produced usable evidence for the HTTP interface, the identification surface, the intervention log, and the application-layer availability, but it cannot verify any property that depends on the cyclic safety communication, which requires the passive-observation or the in-path variant.

\paragraph{Goal achievement.}
The objective was not to certify a product, but to develop and validate a systematic methodology for verifying security functions. This objective was achieved for the central steps of the methodology, which were exercised end to end on a real device. The artifact is best described as a prototype methodology with partial practical validation. It is fully specified across all three network access variants and demonstrated for the tests that the NV1 position supports, while the validation of the NV2 and NV3 tests remains future work. 

\section{Contribution and Practical Relevance}
\label{sec:contribution}

\subsection{Scientific Contributions}
\label{subsec:scientific_contributions}
The work makes five contributions. First, a requirement-derivation model translates the cybersecurity obligations of Regulation (EU) 2023/1230 through prEN 50742 into the component requirements of IEC 62443-4-2, which turns high-level legal obligations into component-level, testable specifications.

Second, a traceability chain links each requirement to an asset, a threat, a security function, a test, and the recorded evidence, so that every executed test serves a distinct requirement and every requirement can be traced to concrete evidence.

Third, the tool-to-target applicability assessment distinguishes an inapplicable tool from a blocked tool and from a genuine failure. This step is necessary, because a general-purpose IT test framework cannot be applied unchanged to an embedded safety gateway, and it prevents a non-execution from being misread as a product defect.

Fourth, the network-access-variant model NV1, NV2, and NV3 makes the feasibility of a test an explicit function of the tester's position on the communication path, which states precisely which properties can and cannot be verified under a given setup.

Fifth, and most importantly, the work demonstrates empirically that a component can enter a safe state after a fault while still failing to satisfy cybersecurity requirements for accountability, integrity, or traceability. The two intervention-log deviations are a concrete example: the device fails safely, yet it does not meet the time-attribution and retention obligations of the regulation. This distinction between a device that fails safely and a device that is secure is the central conceptual result of the thesis. The derivation additionally exposed a gap in the normative basis itself, for example RQ-013, for which prEN 50742 provides no corresponding clause.

\subsection{Industrial Contributions}
\label{subsec:industrial_contributions}
For the industrial partner, the methodology closes the process gap between the regulatory expectations and the established functional test process. It produces the traceable evidence that the technical documentation under the regulation requires, it adapts the existing company tools instead of replacing them, and it yields a reusable schema of requirements, assets, and test cases that can be applied to future product releases and device variants.

\section{Economic and Sustainability Assessment}
\label{sec:economic_sustainability}

\subsection{Economic Perspective}
\label{subsec:economic}
The methodology reduces the effort required for future product releases, because the requirements, assets, and test cases are already linked through a reusable traceability model, so a new release re-runs and extends an existing catalogue rather than deriving tests from scratch. The applicability assessment avoids executing inappropriate tests and prevents engineering effort from being invested in tools that are incompatible with the target device. The structured reuse of the existing company tools limits additional tooling investment. Against these limited costs stands the far larger cost that a security-related field incident or a recall of a safety component would cause, which the early and systematic verification of security functions helps to avoid.

\subsection{Ecological Perspective}
\label{subsec:ecological}
The earlier detection of security weaknesses can extend the safe service life of an installed device and reduce the premature replacement of field hardware, because a weakness that is found and corrected before or during operation avoids a forced device exchange. The reuse of a documented, tool-based method also reduces the repeated setup effort of separate ad-hoc test environments.

\subsection{Social Perspective}
\label{subsec:social}
The verification of cybersecurity functions on safety-related machinery contributes directly to the protection of the humans who operate that machinery. Because the Machinery Regulation treats cybersecurity as a prerequisite for functional safety, a methodology that produces evidence for the security functions of a safety component supports the safety of persons at the machine, which is the primary social objective of the regulation.

\section{Future Work} 
\label{sec:future_work}

\subsection{Immediate Next Steps}
\label{subsec:immediate_next_steps}
The most immediate step is to realise the two remaining network access variants. The passive-observation variant NV2, using a mirroring switch or a network tap, and the in-path variant NV3, using a transparent bridge, would allow the safety-telegram integrity, the in-path manipulation, the watchdog-passivation, and the domain-separation scenarios to be executed, and thereby the requirements that depend on the cyclic safety communication, in particular RQ-002 and RQ-005, to be assessed. The outstanding activities that need no additional variant should be completed in the same campaign: the induced software-installation event for RQ-008, the configuration change for RQ-009, and the TIA-based identification and maintenance reads for RQ-004, RQ-006, RQ-007, and RQ-013. Together these steps close the path from the current partial coverage to a complete requirement assessment. A device that exposes an authenticated and encrypted interface, for example over HTTPS, would additionally enable the credential, transport-security, and access-control tests that were not applicable to the current device.

\subsection{Methodology Extensions}
\label{subsec:methodology_extensions}
Because prEN 50742 grades the required test depth through the Safety-Related Security Level, the catalogue can be extended so that the strictness of a test scales with the level of the affected asset and requirement, for example from a simple checksum check at a low level to cryptographic verification under active manipulation at a high level. The use of machine learning is a further direction. A profiling module that learns the normal network behaviour of a device in the testbed could support the detection of anomalous behaviour during a test \cite[pp. 12]{siboni2016securitytestbedinternetthings}.

\subsection{Organisational Integration}
\label{subsec:organisational_integration}
The methodology should be embedded permanently in the development and quality-management process. This includes assigning the responsibilities of the governance model described in Appendix \ref{app:SecureDevelopmentLifecycle}, fixing the point in the lifecycle at which tests are generated and evaluated, integrating the execution into a continuous integration and delivery pipeline for automated security regression testing, and maintaining the traceability matrix when the normative basis changes, for example when prEN 50742 is published in its final form, so that the methodology can serve as conformity documentation under the Machinery Regulation. A related open point is the precise regulatory scope for a component that is not itself a machine, because the Safe Remote I/O is a component whose obligations arise through its integration into a machine.

%(Security Testbed of source not executed in this thesis: The profiler module for IoT devices is based on a machine learning algorithm which analyzes the network traffic observed in the testbed environment during the testing process. As mentioned in subsection 5.2, the module operates in two phases, a training phase and a testing phase, as illustrated in Figure 11. First, in the training phase, a statistical model is defined based on the network activity recorded during this phase (for that matter, several IoT devices are deployed in the testbed environment). Next, in the testing phase, different IoT devices (that were not used during the
%Figure 11. The IoT profiler module structure, including features extraction, statistical model and decision making submodules.
%training phase) are profiled based on their network behavior, given the statistical model constructed in previous phase \cite[pp. 12]{siboni2016securitytestbedinternetthings})




%TODO für gesamte Arbeit/ Bewertungsschema mit dem diese Bachelorarbeit bewertet wird: Für eine 1.0 müssen folgende Aspekte erfüllt sein
%Fachliche Bearbeitung (Berücksichtigung des Schwierigkeitgrads) --> Sehr gute Abhandlung des Themas bezüglich Lösungswege und Ergebnisse, zusätzliche Alternativen und Aspekte aufgezeigt
%Nutzung von Fachwissen --> Umfassende Kenntnisse, sehr gutes Fachwissen
%Einsatz von Methoden und Werkzeugen --> Bewusste Auswahl, konsequenter und fundierter Einsatz von Ingenieurmethoden und Ingenierwerkzeugen
%Umsetzbarkeit des Ergebnisses --> Sehr gutes Ergebnis ohne Modifikationen umsetzbar oder bereits umgesezt
%Kreativität --> Viele eigene/originelle Ideen eingebracht und in neuen, erfolgreichen Lösungsansätzen umgesetzt
%Wirtschaftliche Bewertung --> Sehr gutes wirtschaftliches Denken in der gesamten Lösung; Zusammenhänge bedacht, nichtfinanzielle Vorteile erkannt
%Berücksichtigung Nachhaltigkeitsaspekte (ökologisch, ökonomisch, sozial) --> Nachhaltigkeit und  Umweltverträglichkeit als zentrales Bewertungskriterium der Arbeit mit fundierter Analyse 
